\scanpage{133}
% archive foliation: 132
Or pour une courbe lisse donnée $X$ sur un schéma de base $S$ de \uline{car. nulle} (p. ex. $S=\operatorname{Spec}\mathbb C$) on a donné une description purement alg. des structures du type a) b) c) (avec connexion $c_X$ sur $P$ rel. à $X/S$), satisfaisant la condition $1^\circ$ --- pour le moment on ne considère pas la condition $2^\circ$. La cat. de ces structures est \uline{rigide}.
De façon précise, les $P$ (fibrés de S.B. de rang 1 sur $X$) sont can. isom. à $P(E)$, où $E$ est l'extension
\diagram{\begin{tikzcd}
0\arrow[r]&\omega_{X/S}\arrow[r]\arrow[d,phantom,"\Vert" sloped]&E\arrow[r]&\mathcal O_X\arrow[r]&0\\
&\underline\omega&&&
\end{tikzcd}}
donnée par
\[E\simeq\mathcal J/\mathcal J^3\otimes\omega^{-1}\]
($\mathcal J\subset\mathcal P^\infty_{X/S}$ l'idéal d'\uline{augmentation}, donc $\mathcal J/\mathcal J^2\simeq\omega_{X/S}=\omega$), la section $\sigma$ est alors définie par cette structure d'extension, enfin la connexion $c$ sur $P=P(E)$ est déduite d'un isom.
\[p:\mathcal P^\infty_{X/S}\xrightarrow{\sim}\mathcal O_{\widehat P}\]
($\widehat P$ complété formel de $P$ le long de $\sigma$), prolongeant l'isom. d'ordre 2 canonique \inserted{$P^{(2)}$} \unsure{impliqué} par la description de $P=P(E)$ par les deux \unsure{premiers} jets de $\mathcal J$. On ne prend, parmi les isom. $p$ \unsure{tels que} la conn. de $\mathcal O_{\widehat P}$ déduite par transport de structure de celle de $\mathcal P^\infty_{X/S}$ (provenant de la conn. tautologique de $X\times_S X$ sur $X$) provienne bien d'\uline{une connexion de $P$}. On constate alors que cet iso (correspondant à
